\begin{table}[tbp]
\centering
\small
\caption{Secondary forward target-selection decomposition. Entries are median text-equal edge-weighted reliability (split IQR) across 100 independent 42/42 reader splits.}
\label{tab:target-selection}
\begin{tabular}{lccc}
\toprule
Specification & Adjacent ($d=1$) & Near skip ($d=2$--$3$) & Far same-line ($d\geq4$) \\
\midrule
Position-only & .9107 [.9081, .9137] & .8533 [.8495, .8570] & .2161 [.2024, .2286] \\
Lexical & .7395 [.7339, .7471] & .6043 [.5918, .6136] & .1392 [.1291, .1530] \\
Syntax & .7346 [.7271, .7416] & .5965 [.5860, .6049] & .1319 [.1244, .1466] \\
Flexible & .7338 [.7257, .7420] & .5959 [.5869, .6064] & .1032 [.0964, .1125] \\
\bottomrule
\end{tabular}
\par\vspace{2pt}\footnotesize Note. Candidate edges were 2,468/4,182/10,119 and eligible edges were 2,444/4,179/10,115 for adjacent/near/far categories. Observed transition mass was 65,789/56,877/3,813. Categories are token separations, not visual-angle saccade amplitudes. See Appendix B for undefined adjacent within-source estimands and permutation resolution.
\end{table}
